\documentclass{article}
\usepackage{icmlko}
\usepackage{graphicx}
\usepackage{booktabs}
\usepackage{amsmath}

\title{무작위로 흩뿌린 사건이 진짜 사건보다 잘 맞았다 \\
  {\large 그리고 위치를 다 뚫어도 상한이 185였다 --- 같은 칸을 두 방향에서 밀고 만난 두 벽}}

\begin{document}
\maketitle

\begin{abstract}
제품 목표 다섯 중 ③「언제 어떤 불특정 이벤트가 발생할지」는 여러 사이클째 \textbf{0} 이다.
이번에는 그 칸 하나를 향해 \textbf{두 팔을 동시에} 보냈다 ---
하나는 격자 8{,}880$\times$577 위에서 \textbf{간격을 직접 재고},
하나는 못 붙던 레코드의 \textbf{위치 키를 뚫는다}.
둘 다 막혔고, \textbf{막힌 이유가 서로 다르다.}
\claimx{조건부 예측이 기후값을 이긴다 --- 개선 $+0.7545$ 일, BCa $[+0.6436, +0.8861]$.}
{구간이 0 을 물면 이 문장은 무너진다. 물지 않았다.}
\claimx{그런데 \textbf{사건 날짜를 격자 안에서 무작위로 섞으면 더 이긴다} --- $+1.0693$, BCa $[+0.9665, +1.1833]$.}
{두 구간이 겹치면 이 문장은 약해진다. 겹치지 않는다 --- 빈틈이 \textbf{0.0804 일}이다.
그리고 이 귀무는 \textbf{진짜 비교보다 먼저} 돌렸다.}
즉 그 이득은 \textbf{시각 구조가 아니라 격자별 발생률}이다(순열이 격자별 사건 수를 보존하므로).
\claimx{다른 팔에서는 벽이 산술이었다 --- 부착을 93 $\to$ 132 로 늘렸지만 \textbf{원리상 상한이 185} 다.}
{팝업 밖 도메인에 위치 열이 하나라도 있으면 상한이 올라간다.
직접 세니 KOBIS 원천 키 5{,}245 중 지리 키 \textbf{0}, 아이돌 키 40 중 \textbf{0} 이었다.}
문턱은 200 이었다. \textbf{지오코딩을 한 건도 하기 전에 계산할 수 있던 수}다.
\end{abstract}

\section{왜 두 팔이었나}

직전 노트는 격자 1.466 억 행을 열고도 \textbf{판 유보 93 행 $\cdot$ 군집 37} 에서 판정 불능이 됐다.
원인은 둘 중 하나여야 했다 --- \textbf{표본이 얇거나}, \textbf{자료에 그 구조가 없거나}.
그래서 이번엔 갈랐다: 팔 ㅌ 은 \textbf{93 행을 아예 안 쓰고} 격자 전체에서 재고,
팔 ㅍ 은 \textbf{93 행 자체를 늘린다}.

\section{팔 ㅌ --- 격자 위에서 간격을 쟀다}

\begin{figure}[h]\centering
\includegraphics[width=\linewidth]{figs/fig1.pdf}
\caption{(a) 순열 귀무와 진짜 비교의 BCa --- \textbf{바닥이 신호보다 위}에 있고 겹치지 않는다.
(b) 팔 넷 중 이긴 것은 하나뿐이다. (c) 문턱 $\tau$ 를 바꿔도 전부 바닥 아래다.}
\end{figure}

사건은 라벨 없이 정의했다 --- 격자별 기준선 대비 $\tau$ 배를 넘는 날.
자는 \textbf{MAE(기후값) $-$ MAE(조건부)}, 군집은 \textbf{격자},
구간은 \texttt{lab.pairboot.cluster\_boot} \textbf{BCa}(B$=$2000)이고 \textbf{폴백 0 건}이다(규약 47).
홀드아웃 간격 \textbf{54{,}862} $\cdot$ 군집 \textbf{2{,}277}.

\subsection{정지 규칙이 실제로 걸렸다}

직전 노트의 설계 실패는 \textbf{문턱을 행 수로 걸고 자의 검출력으로 안 건 것}이었다.
이번엔 사전등록에 검출력 줄을 박았다 --- 사전 추정 MDE $0.04$ 일.
실측 MDE 는 \textbf{0.10838} 로 \textbf{2.7 배 컸다}. 문턱 $0.01\times$기후값MAE $=0.14471$ 아래라 안 멈췄다.

\claimx{\textbf{이 규칙은 장식이 아니었다} --- 실측이 문턱의 75\% 까지 갔다.}
{조금만 더 거칠었으면 진짜 비교를 \emph{계산조차 안 하고} 멈췄을 것이다.
그리고 그 경우에도 이 노트는 나왔을 것이다 --- 「이 표본으로는 못 잰다」가 결론이므로.}

\subsection{판정 --- 이겼는데 졌다}

\begin{table}[h]\centering\small
\begin{tabular}{lrrl}
\toprule
 & 개선(일) & BCa 95\% & 판정 \\
\midrule
🔴 순열 귀무(\textbf{먼저} 돌림) & $+1.0693$ & $[+0.9665, +1.1833]$ & 승 \\
진짜 비교 & $+0.7545$ & $[+0.6436, +0.8861]$ & 승 \\
\midrule
조건부(보조) & $-1.3404$ & $[-1.5193, -1.1811]$ & 패 \\
무정보 & $-0.0574$ & $[-0.0733, -0.0414]$ & 패 \\
\bottomrule
\end{tabular}
\end{table}

\claimx{그러므로 ③ 은 \textbf{오늘도 0} 이다.}
{「lo $>$ 0 이면 승」만 보면 이 노트는 성공으로 읽힌다.
바닥을 같이 재지 않았다면 실제로 그렇게 읽었을 것이다.}
순열은 격자별 \emph{사건 수}를 보존하고 \emph{날짜}만 흩뿌린다.
그런데 흩뿌린 쪽이 더 잘 맞는다 --- 조건부 모형이 얻은 것은
\textbf{「이 격자는 자주 터진다」이지 「이 격자는 지금 터질 참이다」가 아니다.}

\subsection{직전 노트의 「표가 얇아서」는 반만 맞았다}

산출물에서 직접 읽어 계산했다. 직전 이벤트 노트의 BCa 반폭은 기후 MAE 의 \textbf{0.01702\%},
이 노트는 \textbf{0.99976\%} --- \textbf{이 노트의 자가 58.7 배 거칠다}.
그런데 개선율은 $0.00992\% \to 6.221\%$ 로 \textbf{627 배}다.

\claimx{기질을 바꾸면 신호는 정말 커진다. \textbf{바닥이 같이 커져 넘어설 뿐이다.}}
{신호만 커지고 바닥이 안 컸다면 이 노트는 ③ 을 밀었을 것이다.
둘을 같이 재야 이 문장을 쓸 수 있다.}

\subsection{배선 --- 심은 여섯이 전부, 그리고 각각만 발화했다}

\textbf{심은 결함 6 $\cdot$ 발화 6 $\cdot$ 국소 시험 6/6 $\cdot$ 음성 대조 통과 $\cdot$ 검정력 1.0.}
직전 노트의 「죽는 층은 없다」가 \emph{원리상 발화할 수 없는} 검사에서 나왔던 것과 대비된다.
특히 \textbf{W4}(상태 열이 예측에 닿았나)를 심으면 개선이 \textbf{정확히 0.0} 이 되고,
\textbf{W3}(오라클)을 심으면 \textbf{12.58 일}로 튄다 --- 자가 실제로 반응한다.

\section{팔 ㅍ --- 위치 키를 뚫었다. 벽은 산술이었다}

\begin{figure}[h]\centering
\includegraphics[width=\linewidth]{figs/fig2.pdf}
\caption{(a) 부착 전/후와 문턱, 그리고 \textbf{원리상 상한}. (b) 실패를 셋으로 갈라 셈
--- 「지오코딩이 실패했다 43」과 「회복 후에도 좌표 없음 43」은 같은 43 이다.
(c) 교차확인 --- 회복분과 옛 행이 \textbf{같은 자리}에 있다.}
\end{figure}

\subsection{먼저 자기시험을 통과했다}

관을 새로 짜지 않고 직전 노트의 링크 모듈을 \texttt{import} 했고,
회복 전 값이 \textbf{아홉 칸 전부 일치}했다(185/102/83/6/3/93/185/93/67, 격자 37 포함).
\textbf{대조의 출발점이 같다}는 것을 먼저 보인 것이다.

\subsection{회복 --- 그리고 상한}

\begin{table}[h]\centering\small
\begin{tabular}{lrrrl}
\toprule
 & 전 & 후 & $\Delta$ & 문턱 \\
\midrule
붙는 행 & 93 & \textbf{132} & $+39$ & $\ge$200 ✗ \\
서로 다른 격자(군집) & 37 & \textbf{42} & $+5$ & $\ge$60 ✗ \\
\bottomrule
\end{tabular}
\end{table}

\claimx{\textbf{이 자를 못 넘었다} --- 그래서 프로브를 \textbf{안 돌렸다}.}
{소수 라벨 곡선도 판 $\rho$ 도 안 쟀다. 「그래도 돌려는 봤다」를 안 했다.}
보수적으로 세도 판정이 안 바뀐다: 격자 모호분을 버리면 108/41, 교차확인 어긋남까지 버리면 103/38.

그리고 벽의 정체는 산술이다. 부착의 모수는 \textbf{팝업 두 도메인의 유보 합집합 185} 이고
\textbf{문턱은 200} 이었다. 팝업 밖에 위치 열이 있으면 상한이 오르는데, 직접 세니
\textbf{KOBIS 원천 키 5{,}245 중 지리 키 0} $\cdot$ \textbf{아이돌 키 40 중 0} 이고
축에 걸린 것은 배급사 저명도 스칼라뿐이었다.

\claimx{\textbf{이 계산은 지오코딩 한 건 전에 할 수 있었다.}}
{그리고 실제로 그렇게 했다 --- 사전등록의 검출력 줄에 「원리상 상한 185 $<$ 문턱 200」을 미리 적었다.
직전 노트가 이 칸이 없어서 원리상 못 가르는 자리까지 갔던 바로 그 자리다.}

\subsection{어긋남은 회복이 만든 것이 아니다}

좌표의 구가 레코드 자신이 적은 구와 같은지 세었다 ---
회복분 \textbf{0.641}(25/39), \textbf{회복 전부터 붙던 행 0.6374}(58/91).

\claimx{거의 같다 --- \textbf{어긋남은 이 자료가 원래 갖고 있던 성질}이다.}
{옛 행에 같은 자를 안 댔다면 「회복분이 나쁘다」로 읽혔을 것이다.
음성 대조가 없으면 「나쁘다」를 말할 기준 자체가 없다.}

\subsection{HTTP 200 이 두 번 거짓말했다}

무료 공공 지오코더 둘에 키 없이 요청했다(2 건 $\cdot$ 광역 크롤 아님).
\textbf{둘 다 HTTP 200} 이고 본문이 각각
\texttt{E0001 승인되지 않은 KEY} $\cdot$ \texttt{PARAM\_REQUIRED key가 없어서} 였다.

\claimx{그러므로 이것은 「무료 지오코딩이 없다」가 아니라 \textbf{「키가 필요해서 못 했다」}이다.}
{응답 원문을 산출물에 통째로 실었다. 상태 코드만 봤다면 성공 2 건으로 회계됐을 것이다.}

\section{자기 적발 --- 같은 코드가 132 와 129 를 냈다}

\claimx{집합 순회 순서가 프로세스마다 달라 좌표가 250m 경계를 넘나들었다.}
{\textbf{종료 코드는 계속 0} 이었다. 산출물도 정상으로 보였다.
결정성 시험을 신설하고 3 회 연속 동일을 확인한 뒤에야 위 수를 썼다.}

\section{한계 --- 「없다」가 아니라 「못 했다」}

\begin{itemize}
\item \textbf{순열이 발생률은 안 부순다} --- 발생률까지 부수는 귀무를 \textbf{안 돌렸다}.
      「시각 구조가 없다」는 이 노트가 쓴 문장의 상한이다
\item 행 가중과 격자 가중의 \textbf{부호가 다른데} 어느 쪽이 자인지 미리 안 적어 판정에 안 썼다
\item 민감도 셋에는 귀무를 \textbf{안 돌렸다} --- 그림 (c) 의 바닥선은 $\tau=1.2$ 의 것이다
\item 관측률이 낮은 격자 \textbf{521 개}를 뺐다
\item ③ 의 \textbf{「어떤」은 안 밀었다} --- 이 노트가 민 것은 「언제」뿐이다
\item 🔴 시운전 뒤 러너에 \textbf{진단 절과 직렬화 청소를 더했다.}
      팔$\cdot$문턱$\cdot$판정 규칙은 안 건드렸으나 \textbf{그 기계 증거는 없다}
\item 동/구 단위로는 22 건이 더 붙지만 \textbf{격자 부착에 안 더했다} --- 다른 분모다
\item 사전등록이 산출물과 같은 커밋에 들어가므로 \textbf{순서의 증거는 mtime 뿐}이다.
      두 팔 모두 그 약점을 \textbf{사전등록 안에 스스로 먼저 적었다}
\end{itemize}

\end{document}
